\documentclass[11pt,a4paper]{article}
\usepackage[T1]{fontenc}
\usepackage[utf8]{inputenc}
\usepackage{lmodern,microtype}
\usepackage[margin=27mm]{geometry}
\usepackage{amsmath,amssymb,amsthm,mathtools}
\usepackage{booktabs,array,longtable,enumitem,xcolor}
\usepackage[colorlinks=true,linkcolor=blue!45!black,citecolor=blue!45!black,urlcolor=blue!45!black]{hyperref}
\usepackage{fancyhdr}
\pagestyle{fancy}\fancyhf{}
\fancyhead[L]{\small Frobenius lifting, rigidity, and stabilization}
\fancyhead[R]{\small T76--T80}\fancyfoot[C]{\thepage}
\setlength{\headheight}{14pt}
\setlength{\parskip}{4pt}\setlength{\parindent}{1.1em}
\setlist{itemsep=3pt,topsep=5pt}
\newtheorem{theorem}{Theorem}
\renewcommand{\thetheorem}{T\arabic{theorem}}
\setcounter{theorem}{75}
\newtheorem{lemma}{Lemma}[section]
\newtheorem{proposition}[lemma]{Proposition}
\newtheorem{corollary}[lemma]{Corollary}
\theoremstyle{definition}
\newtheorem{definition}[lemma]{Definition}
\newtheorem{problem}[lemma]{Auxiliary problem}
\theoremstyle{remark}\newtheorem{remark}[lemma]{Remark}
\newcommand{\Z}{\mathbb Z}\newcommand{\Q}{\mathbb Q}
\newcommand{\C}{\mathbb C}\newcommand{\F}{\mathbb F}
\DeclareMathOperator{\rank}{rank}
\hypersetup{pdftitle={Frobenius Lifting, Floquet Rigidity, and Equivariant Stabilization: T76-T80},pdfauthor={KEIO}}
\title{Frobenius Lifting, Floquet Rigidity,\\and Equivariant Stabilization\\[5pt]\large Five theorems and their applications to prior open problems}
\author{KEIO}
\date{24 October 2026\\[3pt]\small Research manuscript T76--T80}
\begin{document}
\maketitle
\begin{abstract}
We develop explicit constructions in three settings. For the lacunary series
$\Phi_p(z)=\sum_{n\geq0}z^{p^n}$, a Frobenius-compatible correction tower
produces a monic polynomial basis with exact evaluation valuations $v_p(d!)$.
We determine the full kernel of evaluation modulo every prime power and prove
the minimum-degree conjecture of Krattenthaler and M\"uller, first posed in
2015. A multiplicative version, followed by a finite descent of denominators,
proves their 2013 minimum-degree conjecture for the Cantor product
$\prod_{n\geq0}(1+z^{3^n})$ in its original coefficient ring.
For discrete Schr\"odinger operators, coordinate-average descent and the
second and third Fourier perturbation identities prove rigidity at period
$(3\mathbb Z)^d$ in every dimension. Combined with prior results, this completes
the classification for rectangular period lattices.
Finally, for an order-two lattice action, we construct mixed linear and
quadratic systems of parameters with an optimal quadratic budget. Their
Artinian quotients prove Clarke's coefficient inequality and give the sharp
universal stabilization exponent $\max(p-q,0)$.
These results comprise three numbered-conjecture applications and one
explicitly posed, unnumbered classification problem. The five theorem labels
do not denote five independently posed open problems.
\end{abstract}

\section{Research questions, exact targets, and scope}

The organizing strategy is to pose an auxiliary construction problem, solve
it with explicit control of the quantity that obstructs a prior problem,
and then return to that original problem without changing its hypotheses.
The auxiliary quantities here are evaluation valuation, low-order Fourier
constraints, and the number of quadratic parameters.

\begin{center}\small
\begin{tabular}{@{}p{.08\textwidth}p{.43\textwidth}p{.40\textwidth}@{}}
\toprule
Label & Structural result & External application \\
\midrule
T76 & A factorial-valued monic basis, exact valuation formula, and the full evaluation kernel. & Auxiliary construction used in T77 and adapted in T78.\\
T77 & Degree-optimal integral annihilators for every prime power. & Krattenthaler--M\"uller Conjecture~4, first preprint 2015 \cite{KMp}.\\
T78 & Multiplicative lifting with finite coefficient-ring descent. & Krattenthaler--M\"uller Conjecture~12, 2013 \cite{KM3}.\\
T79 & All-dimensional period-three rigidity and the rectangular classification. & The remaining small-period problem in \cite[Section~4.4]{CFK}.\\
T80 & Optimal mixed-degree parameter systems and sharp universal stabilization. & Clarke Conjecture~8.5, 2026 \cite{Clarke26}.\\
\bottomrule
\end{tabular}
\end{center}

The two arithmetic targets concern formal series modulo integer powers, not
the arithmetic nature of a numerical value. The spectral target is the
rectangular-period problem, not the larger classification over arbitrary
oblique lattices. The Ehrhart result applies without assuming that the
equivariant numerator is already polynomial or effective. These distinctions
are retained throughout the proofs.

The claims below are unconditional mathematical statements with complete
arguments, using the established results explicitly cited. In particular,
the Cantor lower bound is the prior Theorem~10 of \cite{KM3}; the remaining
rectangular cases use the prior period-reduction and existence theorems of
\cite{CFK}; and the Ehrhart argument uses the Cohen--Macaulay property of
normal affine semigroup rings. Exact finite computations supplement these
arguments and are not used to extrapolate an all-parameter claim.

This is a research draft, independently checked internally but not externally
peer reviewed or proof-assistant certified. The current literature search
located no earlier identical resolution of the stated targets. The
Frobenius congruence, Fourier perturbation, pole extraction, generic sections,
and Koszul complexes are established mathematics. We therefore do not assert
that the underlying general principles are historically unprecedented.
The attribution and success-condition audit at the end makes this limitation
explicit.


\section{A Frobenius-compatible lifting problem}
\label{arith:section}

The first two existing targets concern sharp algebraic relations modulo prime
powers for transcendental formal series. Krattenthaler and M\"uller
\cite[Conjecture~4]{KMp} ask for the minimum degree of a monic polynomial
annihilating
\[
 \Phi_p(z)=\sum_{n\geq0}z^{p^n}
\]
coefficientwise modulo $p^\gamma$, with coefficients in $\mathbb Z[z]$.
Their proposed answer is the least $d$ satisfying $\nu_p(d!)\geq\gamma$.
Here $p$ is an arbitrary prime and $\gamma\geq1$.
The binary version is also present in \cite[Conjecture~4]{KKM2}; it is a
special case of this target, not an additional independent problem.
The second target, \cite[Conjecture~12]{KM3}, concerns the Cantor product
\[
 \Psi(z)=\prod_{j\geq0}(1+z^{3^j}).
\]
Its prescribed coefficient ring is
$\mathbb Z[z,z^{-1},(1+z)^{-1}]$, and the proposed minimum degree modulo
$3^\gamma$ is $2d$, for the least positive $d$ with
$d+\nu_3(d!)\geq\gamma$.

We begin with an auxiliary question that provides the same constructive
step for both targets: given an integral series with an affine functional
equation under $z\mapsto z^p$, can a monic relation of known divisibility
be lifted recursively while controlling its degree? The answer uses the
classical Frobenius congruence
\begin{equation}\label{arith:frobenius}
 B(z)^p\equiv B(z^p)\pmod p\qquad(B\in\mathbb Z[[z]]).
\end{equation}
The resulting substitution-dependent recursion supplies the correction
polynomials explicitly. The underlying congruence is classical; its use here
is not described as a historically new foundational principle.

For an integer polynomial or integral formal series $G$, write
$\nu_p(G)$ for the minimum $p$-adic valuation of its coefficients, with
$\nu_p(0)=+\infty$. Congruence of formal series always means coefficientwise
congruence.

\begin{lemma}[Lifting through an affine Frobenius equation]
\label{arith:lift}
Let $R\subseteq\mathbb Z[[z]]$ be a ring stable under the substitution
$\sigma(c)(z)=c(z^p)$. Suppose $F\in\mathbb Z[[z]]$ satisfies
\[
 F(z^p)=a(z)F(z)+b(z),\qquad a,b\in R.
\]
For $G\in R[t]$, define
\[
 (\tau G)(z,t)=\sum_j \sigma(c_j)(z)(a(z)t+b(z))^j
 \quad\text{if }G=\sum_jc_j(z)t^j.
\]
For a nonnegative integer $e$, if $G$ is monic of degree $D>0$ and
$G(z,F(z))\in p^e\mathbb Z[[z]]$, then
\begin{equation}\label{arith:liftoperator}
 \mathcal L_e(G)=G^p-p^{(p-1)e}\tau G
\end{equation}
is monic of degree $pD$ and has evaluation divisible by $p^{pe+1}$.

Starting with a monic $G_0$ of degree $D_0>0$ and evaluation divisible
by $p^{e_0}$ for a nonnegative integer $e_0$, set
\begin{equation}\label{arith:generalrecurrence}
 G_{r+1}=G_r^p-p^{(p-1)e_r}\tau G_r,
 \qquad e_{r+1}=pe_r+1.
\end{equation}
Then $G_r$ is monic of degree $D_0p^r$ and has evaluation divisible by
$p^{e_r}$, where
\[
 e_r=p^re_0+\frac{p^r-1}{p-1}.
\]
If $d=\sum_r d_rp^r$ is a base-$p$ expansion, the product
$\prod_rG_r^{d_r}$ is monic of degree $D_0d$ and has evaluation divisible
by $p^{e_0d+\nu_p(d!)}$.
\end{lemma}
\begin{proof}
Write $G(z,F(z))=p^eB(z)$ with $B\in\mathbb Z[[z]]$. The affine functional
equation makes evaluation commute with the two substitutions:
\[
 (\tau G)(z,F(z))=G(z^p,F(z^p))=p^eB(z^p).
\]
Consequently
\[
 \mathcal L_e(G)(z,F(z))
   =p^{pe}\bigl(B(z)^p-B(z^p)\bigr),
\]
which is divisible by $p^{pe+1}$ by \eqref{arith:frobenius}.
The correction $\tau G$ has $t$-degree at most $D$, whereas $G^p$ is
monic of degree $pD>D$. This proves the first assertion and permits
induction in \eqref{arith:generalrecurrence}.
The recurrence for $e_r$ has the displayed solution. Finally, Legendre's
formula gives
\[
 \sum_r d_re_r=e_0d+\sum_r d_r\frac{p^r-1}{p-1}
       =e_0d+\frac{d-s_p(d)}{p-1}=e_0d+\nu_p(d!),
\]
where $s_p(d)=\sum_rd_r$. Multiplying the evaluation divisibilities
proves the last assertion.
\end{proof}

\section{An exact factorial basis and its evaluation kernel}

For $\Phi_p$, the exact identity
\begin{equation}\label{arith:phifunctional}
 \Phi_p(z^p)=\Phi_p(z)-z
\end{equation}
allows us to take $R=\mathbb Z[z]$, $a=1$, $b=-z$, and the seed $A_0=t$.
Put $e_r=(p^r-1)/(p-1)$ and define
\begin{equation}\label{arith:Ar}
 A_0(z,t)=t,\qquad
 A_{r+1}(z,t)=A_r(z,t)^p-p^{(p-1)e_r}A_r(z^p,t-z).
\end{equation}
All substitutions in this formula are simultaneous. For $d\geq0$ with
base-$p$ digits $d_r$, set
\begin{equation}\label{arith:Bd}
 B_d(z,t)=\prod_{r\geq0} A_r(z,t)^{d_r};\qquad B_0=1.
\end{equation}
Only finitely many factors occur. Lemma~\ref{arith:lift} implies that $B_d$
is monic of degree $d$ and that
\begin{equation}\label{arith:upperval}
 \nu_p\bigl(B_d(z,\Phi_p(z))\bigr)\geq\nu_p(d!).
\end{equation}

The opposite bound will follow from an explicit coefficient-isolation
argument. In particular, it does not require a linear-independence theorem
for a separate family of auxiliary series.

\begin{lemma}[Isolation by digit sums]\label{arith:isolate}
Let $H(z,t)=\sum_{k=0}^{D}c_k(z)t^k\in\mathbb Z[z,t]$, where $D\geq1$,
and choose $M\geq0$ with $\deg c_k\leq M$ for every nonzero $c_k$.
Choose $L\geq0$ such that $p^L>M$, set $K=L+2$, and define
\[
 N=\sum_{i=0}^{D-1}p^{K+i}.
\]
Then, for every $j_0\in\{0,\ldots,M\}$,
\begin{equation}\label{arith:isolationidentity}
 [z^{N+j_0}]H(z,\Phi_p(z))=D!\,[z^{j_0}]c_D(z).
\end{equation}
In particular, if $H(z,\Phi_p(z))\equiv0\pmod{p^\gamma}$, then every
coefficient of $D!c_D(z)$ is divisible by $p^\gamma$.
\end{lemma}
\begin{proof}
Every exponent appearing in $\Phi_p(z)^k$ is a sum of $k$ powers of $p$.
Its base-$p$ digit sum is at most $k$, since every carry decreases the
number of summands by $p-1$.

Consider the term $[z^j]c_k(z)$ in a coefficient extraction at $N+j_0$,
and put $\Delta=j_0-j$. If $\Delta>0$, then $\Delta\leq M<p^K$ and
\[
 s_p(N+\Delta)=D+s_p(\Delta)>D.
\]
If $\Delta=-c<0$, then $1\leq c\leq M<p^L$ and borrowing at the
lowest nonzero digit of $N$ gives
\begin{align*}
 s_p(N-c)
   &=D-1+s_p(p^K-c)\\
   &=D-1+(p-1)K-s_p(c-1)\\
   &\geq D-1+(p-1)(K-L)>D.
\end{align*}
Here $s_p(c-1)\leq(p-1)L$, and $K-L=2$. Thus, whenever $j\ne j_0$,
the required coefficient in every $\Phi_p^k$, $k\leq D$, is zero.

For $j=j_0$, the exponent is $N$ and has digit sum $D$. Therefore only
$k=D$ can contribute. A representation of $N$ by exactly $D$ powers of
$p$ cannot involve a carry: a carry would lower the final digit sum below
$D$. Its multiset of summands is therefore precisely the $D$ distinct
powers in the definition of $N$. Their ordered arrangements number $D!$.
This proves \eqref{arith:isolationidentity} and the divisibility assertion.
\end{proof}

For later statements, define
\[
 D_\gamma=\min\{d\geq0:\nu_p(d!)\geq\gamma\},\qquad
 J_\gamma=\{H\in\mathbb Z[z,t]:H(z,\Phi_p(z))\equiv0\pmod{p^\gamma}\}.
\]
The set $J_\gamma$ is an ideal because evaluation modulo $p^\gamma$ is
a ring homomorphism.

\begin{theorem}[Factorial basis, exact valuation, and the full kernel]
\label{arith:factorialbasis}
The polynomials $B_d$ in \eqref{arith:Bd} form a monic triangular
$\mathbb Z[z]$-basis of $\mathbb Z[z,t]$. For every finite expansion
$H=\sum_dc_d(z)B_d(z,t)$,
\begin{equation}\label{arith:exactnorm}
 \nu_p\bigl(H(z,\Phi_p(z))\bigr)
   =\min_d\bigl(\nu_p(c_d)+\nu_p(d!)\bigr).
\end{equation}
In particular, each $B_d$ has evaluation valuation exactly $\nu_p(d!)$.
For every $\gamma\geq1$, membership in $J_\gamma$ is equivalent to
\begin{equation}\label{arith:kernelcriterion}
 c_d(z)\in p^{\max\{0,\gamma-\nu_p(d!)\}}\mathbb Z[z]
 \quad\text{for every }d.
\end{equation}
The ideal has the finite presentation
\begin{equation}\label{arith:kernelideal}
 J_\gamma=
 \left(B_{D_\gamma},\;
 p^{\gamma-\nu_p(d!)}B_d\ :\ 0\leq d<D_\gamma\right),
\end{equation}
and, as a $\mathbb Z[z]$-module,
\begin{equation}\label{arith:module}
 \mathbb Z[z,t]/J_\gamma
 \simeq\bigoplus_{d=0}^{D_\gamma-1}
       \bigl(\mathbb Z/p^{\gamma-\nu_p(d!)}\mathbb Z\bigr)[z].
\end{equation}
Moreover, giving $t$ weight $1$ and $z$ weight $p$ yields
$\operatorname{wt}_p(B_d)\leq d$; consequently
$\deg_z B_d\leq\lfloor d/p\rfloor$.
\end{theorem}
\begin{proof}
Since $B_d$ is monic of $t$-degree $d$, successively subtracting its
leading term gives a unique finite expansion of every polynomial in the
$B_d$. This proves the basis assertion.

Sufficiency of \eqref{arith:kernelcriterion} follows from
\eqref{arith:upperval}. For necessity, suppose $H\in J_\gamma$ and let
$D$ be the largest index with $c_D\ne0$. This is also its ordinary
$t$-degree, with leading coefficient $c_D(z)$. If $D\geq1$,
Lemma~\ref{arith:isolate} shows
\[
 c_D(z)\in p^{\max\{0,\gamma-\nu_p(D!)\}}\mathbb Z[z].
\]
The same conclusion is immediate for $D=0$. By
\eqref{arith:upperval}, the evaluation of $c_DB_D$ vanishes modulo
$p^\gamma$. Subtracting it leaves a smaller-degree element of $J_\gamma$.
Induction proves \eqref{arith:kernelcriterion} for every coefficient.

For $H\ne0$, put $a=\min_d(\nu_p(c_d)+\nu_p(d!))$, a finite
nonnegative integer. Equation~\eqref{arith:upperval} gives valuation at
least $a$. If the valuation were at least $a+1$, criterion
\eqref{arith:kernelcriterion} with $\gamma=a+1$ would contradict the
definition of $a$. This proves \eqref{arith:exactnorm}; the zero case
uses the convention $\nu_p(0)=+\infty$.

Each displayed generator in \eqref{arith:kernelideal} belongs to
$J_\gamma$. Conversely, divide any $H\in J_\gamma$ by the monic
$B_{D_\gamma}$ in $\mathbb Z[z][t]$. The remainder is in $J_\gamma$
and has degree below $D_\gamma$. Its unique basis expansion satisfies
\eqref{arith:kernelcriterion}, placing it in the displayed ideal.
Uniqueness of the remainder coefficients modulo their indicated powers
of $p$ also proves the module decomposition \eqref{arith:module}.
This is a module isomorphism; no diagonal description of multiplication
in the quotient ring is asserted.

Finally, $A_0=t$ has weight one. The substitution $z\mapsto z^p$,
$t\mapsto t-z$ multiplies the weight of every monomial by at most $p$:
the substituted $z$ has weight $p^2$ and the substituted $t$ has weight
at most $p$. Thus \eqref{arith:Ar} gives
$\operatorname{wt}_p(A_r)\leq p^r$ inductively. Multiplying according
to the base-$p$ digits of $d$ yields
$\operatorname{wt}_p(B_d)\leq\sum_rd_rp^r=d$, proving the last assertion.
\end{proof}

The coefficient-isolation identity explains why the upper bound from
Frobenius lifting is sharp: it produces an exact coefficient $d!$ at a
specified exponent for every monic polynomial of degree $d$ once its
coefficient degrees have been bounded. It also prevents cancellation
between different valuation levels in the basis expansion.

\section{The original minimum-degree conjecture for \texorpdfstring{$\Phi_p$}{Phi p}}

\begin{theorem}[Krattenthaler--M\"uller Conjecture~4]
\label{arith:phiconjecture}
For every prime $p$ and every integer $\gamma\geq1$, the smallest
$t$-degree of a monic polynomial $H(z,t)\in\mathbb Z[z,t]$ satisfying
\[
 H(z,\Phi_p(z))\equiv0\pmod{p^\gamma}
\]
is precisely
\[
 D_\gamma=\min\{d\geq0:\nu_p(d!)\geq\gamma\}.
\]
A degree-optimal polynomial is $B_{D_\gamma}$, defined explicitly by
\eqref{arith:Ar}--\eqref{arith:Bd}; it has
$z$-degree at most $\lfloor D_\gamma/p\rfloor$.
This proves the full conjecture in \cite{KMp}, with its original
coefficient ring and quantifiers.
\end{theorem}
\begin{proof}
Equation~\eqref{arith:upperval} shows that $B_{D_\gamma}$ is a monic
annihilator modulo $p^\gamma$ of degree $D_\gamma$. Conversely, for a
monic annihilator of degree $D\geq1$, Lemma~\ref{arith:isolate} gives
$p^\gamma\mid D!$. Hence $D\geq D_\gamma$. Degree zero cannot occur:
the only monic constant polynomial is one. The coefficient-polynomial degree bound follows
from Theorem~\ref{arith:factorialbasis}.
\end{proof}

For illustration, when $p=2$, the recursion begins with
\begin{align*}
 A_1(z,t)&=t^2-t+z,\\
 A_2(z,t)&=t^4-2t^3+(2z-1)t^2+(2z+2)t-3z^2-2z.
\end{align*}
The evaluations of these polynomials have valuations exactly one and
three. Thus $A_2$ is a minimum-degree relation modulo $8$. At the next
levels, $A_3$ and $A_4$ have degrees eight and sixteen and evaluation
valuations seven and fifteen. These statements hold for the full formal
series, not only a truncation. The recurrence supplies the polynomials
without solving an undetermined coefficient system at each level.

\section{The original Cantor-product minimum-degree conjecture}

For this second application, the affine functional equation is
\begin{equation}\label{arith:psifunctional}
 \Psi(z^3)=\frac{\Psi(z)}{1+z}.
\end{equation}
The original coefficient ring
$R=\mathbb Z[z,z^{-1},(1+z)^{-1}]$ is not stable under substitution
$z\mapsto z^3$ as an ordinary rational-function ring. We first perform
the construction in the stable ring
\[
 R_\infty=\mathbb Z\bigl[z,(1+z^{3^j})^{-1}:j\geq0\bigr]
       \subseteq\mathbb Z[[z]],
\]
and then descend every coefficient at the required finite precision.
The following lemma ensures that the final answer uses the prescribed
ring, rather than a changed problem.

\begin{lemma}[Finite-precision descent]\label{arith:descent}
For every $\gamma\geq1$ and every $c\in R_\infty$, there exists
$\widetilde c\in\mathbb Z[z,(1+z)^{-1}]$ such that
$c-\widetilde c\in3^\gamma\mathbb Z[[z]]$.
Consequently a monic polynomial in $R_\infty[t]$ can be replaced, without
changing its degree or any evaluation modulo $3^\gamma$ at an integral
formal series, by a monic polynomial in
$\mathbb Z[z,(1+z)^{-1}][t]$.
\end{lemma}
\begin{proof}
For $j\geq0$ put $U_j=1+z^{3^j}$ and $V_j=(1+z)^{3^j}$.
Frobenius modulo three gives $U_j=V_j+3E_j$ for some $E_j\in\mathbb Z[z]$.
The finite geometric identity gives
\begin{equation}\label{arith:finiteinverse}
 U_j^{-1}\equiv
 \sum_{a=0}^{\gamma-1}(-3E_j)^aV_j^{-a-1}
 \pmod{3^\gamma\mathbb Z[[z]]}.
\end{equation}
All denominators on the right are powers of $1+z$. More explicitly,
the omitted remainder is $(-3E_j)^\gamma V_j^{-\gamma}U_j^{-1}$,
whose formal coefficients are divisible by $3^\gamma$ because both
$U_j$ and $V_j$ have constant term one.
Every $c\in R_\infty$ is a polynomial expression in $z$ and finitely many
$U_j^{-1}$. Substituting \eqref{arith:finiteinverse} in that expression
proves the first assertion. For a polynomial in $t$, perform the
replacement on its finitely many non-leading coefficients and keep the
leading coefficient one unchanged. Evaluation at an integral series
preserves the congruence.
\end{proof}

Set
\begin{equation}\label{arith:Cr}
 \begin{split}
 C_0(z,t)&=t^2-\frac1{1+z},\qquad f_0=1,\\
 C_{r+1}(z,t)&=C_r(z,t)^3-
       3^{2f_r}C_r\left(z^3,\frac{t}{1+z}\right),
       \qquad f_{r+1}=3f_r+1.
 \end{split}
\end{equation}
The seed has evaluation divisible by three. Indeed,
$\Psi(z)^3\equiv\Psi(z^3)\pmod3$ by Frobenius, and
\eqref{arith:psifunctional} implies
$\Psi(z)^2\equiv(1+z)^{-1}\pmod3$ because $\Psi$ has constant term one
and is a unit of $\mathbb Z[[z]]$.
Applying Lemma~\ref{arith:lift} inside $R_\infty$ gives
\begin{equation}\label{arith:Cproperties}
 \deg_t C_r=2\cdot3^r,\qquad
 f_r=\frac{3^{r+1}-1}{2},\qquad
 C_r(z,\Psi(z))\in3^{f_r}\mathbb Z[[z]],
\end{equation}
with every $C_r$ monic.

\begin{theorem}[Krattenthaler--M\"uller Conjecture~12]
\label{arith:cantorconjecture}
Let $\gamma\geq1$ and let
\[
 d_\gamma=\min\{d\geq1:d+\nu_3(d!)\geq\gamma\}.
\]
The smallest degree of a monic polynomial
$H(z,t)\in\mathbb Z[z,z^{-1},(1+z)^{-1}][t]$ satisfying
\[
 H(z,\Psi(z))\equiv0\pmod{3^\gamma}
\]
is exactly $2d_\gamma$. A polynomial attaining this degree is obtained
from the base-three digit product of \eqref{arith:Cr} followed by the
finite descent in Lemma~\ref{arith:descent}. In particular, the unchanged
Conjecture~12 of \cite{KM3} holds for every $\gamma$.
\end{theorem}
\begin{proof}
Write $d=d_\gamma=\sum_r d_r3^r$, where $0\leq d_r<3$, and put
$E_d=\prod_r C_r^{d_r}\in R_\infty[t]$. It is monic of degree $2d$.
Using \eqref{arith:Cproperties} and Legendre's formula,
\[
 \nu_3\bigl(E_d(z,\Psi(z))\bigr)
  \geq\sum_rd_rf_r
  =\frac{3d-s_3(d)}2
  =d+\nu_3(d!)\geq\gamma.
\]
Lemma~\ref{arith:descent} replaces $E_d$ by a monic polynomial
$\widetilde E_d$ of the same degree in
$\mathbb Z[z,(1+z)^{-1}][t]$, with the same evaluation modulo
$3^\gamma$. This coefficient ring is a subring of the ring prescribed
in the original problem, so it supplies the required upper bound
without enlarging the allowed ring.

For the reverse inequality, Theorem~10 of \cite{KM3} proves that the
minimum degree in precisely that prescribed ring is even, say $2b$,
and must satisfy $b+\nu_3(b!)\geq\gamma$. Hence $b\geq d_\gamma$.
Combined with the constructed degree-$2d_\gamma$ polynomial, this
proves equality. The cited lower bound is an established theorem;
the conjectural step in that source was attainment of the bound for
all $\gamma$, which the construction above provides.
\end{proof}

An explicit example in the original ring is the degree-six polynomial
\[
 t^6-\frac{3t^4}{1+z}-\frac{6t^2}{(1+z)^2}
       +\frac8{(1+z)^3}+\frac{27z}{(1+z)^5},
\]
whose evaluation at $\Psi$ vanishes modulo $81$. To obtain it from
$C_1$, use
\[
 \frac1{1+z^3}\equiv\frac1{(1+z)^3}
          +\frac{3z}{(1+z)^5}\pmod9.
\]
Only precision nine is needed for this coefficient because it is
multiplied by nine in $C_1$. Since $3+\nu_3(3!)=4$, the theorem
certifies that degree six is minimal for modulus $3^4$.

The two applications share the seed-lifting mechanism, but their sharpness
arguments and coefficient rings differ. For $\Phi_p$ the coefficient
isolation also determines the entire evaluation kernel and its module
structure. For $\Psi$, finite-precision descent is essential to preserve
the original problem's ring. The binary specialization of the first
application remains nested within the all-prime result.

\section{Small-period Floquet rigidity in every dimension}
\label{sec:ternary-rigidity}

Let $\Delta$ be the discrete Laplacian on $\Z^d$ with nearest-neighbor
coefficients equal to one. Two complex periodic potentials are called
\emph{Floquet isospectral} when their Floquet matrices have identical
characteristic polynomials at every quasimomentum. We address the remaining
rectangular-period problem explicitly identified in
\cite[Section~4.4]{CFK}: whether a nonzero complex potential of period
$(3\Z)^d$, $d\geq3$, can be Floquet isospectral to zero. That source singles
out $d=3$ as the smallest unresolved case. This is an explicitly posed but
unnumbered question. It is narrower than the classification over arbitrary
oblique lattices, which is not settled here.

Our auxiliary question is whether coordinate averaging, followed by just the
second and third perturbation identities, can force the potential to vanish.
The proof below gives an affirmative answer. Fourier representations, rational
spectral identities, and extraction of coefficients through their pole
hyperplanes have prior applications to Floquet rigidity; see, in particular,
\cite[Lemma~3.2, Theorem~3.3, and Section~4]{Liu23}. The contribution proposed
here is their combination with coordinate-average descent and the specific
third-order constraint for the period-three Fourier support.

\begin{theorem}[Period-three rigidity and the rectangular classification]
\label{thm:ternary-rigidity}
For every $d\geq1$, a complex $(3\Z)^d$-periodic potential $V$ on $\Z^d$
that is Floquet isospectral to zero satisfies $V\equiv0$.
Consequently, for arbitrary positive integers $q_1,\ldots,q_d$, there exists
a nonzero complex potential of period
$q_1\Z\oplus\cdots\oplus q_d\Z$ Floquet isospectral to zero if and only if
\begin{equation}\label{eq:rectangular-threshold}
 \max_{1\leq j\leq d}q_j\geq4.
\end{equation}
\end{theorem}

We prove the period-three assertion directly. The final classification then
uses the previously established existence and period-reduction results of
\cite{CFK}.

\subsection{The Floquet gauge and coordinate-average descent}

Write $G_d=(\Z/3\Z)^d$ and $\omega=e^{2\pi i/3}$. For
$t=(t_1,\ldots,t_d)\in(\C^*)^d$, use the Floquet gauge
$f(n)=t_1^{n_1}\cdots t_d^{n_d}\phi(n)$, where $\phi$ is three-periodic
in each coordinate. The Floquet multipliers are $t_j^3$, and the gauged
matrix on functions on $G_d$ is
\begin{equation}\label{eq:ternary-gauge}
 H_V(t)=\sum_{j=1}^d(t_jS_j+t_j^{-1}S_j^{-1})+\operatorname{diag}(V),
 \qquad (S_j\phi)(n)=\phi(n+e_j).
\end{equation}
Floquet isospectrality gives
\begin{equation}\label{eq:ternary-charidentity}
 \det(\lambda I-H_V(t))=\det(\lambda I-H_0(t))
 \qquad(t\in(\C^*)^d).
\end{equation}
Indeed, the equality initially holds on the product unit torus. Each
characteristic-coefficient difference is a Laurent polynomial in $t$;
the one-variable polynomial identity theorem, applied successively to its
variables, extends the equality to $(\C^*)^d$.

For a coordinate $j$, define the averaged potential on $G_{d-1}$ by
\[
 (\mathcal A_jV)(n_1,\ldots,\widehat{n_j},\ldots,n_d)
   =\frac13\sum_{r=0}^2V(n_1,\ldots,r,\ldots,n_d).
\]

\begin{lemma}[Coordinate-average descent]\label{lem:average-descent}
If $V$ satisfies \eqref{eq:ternary-charidentity}, then $\mathcal A_jV$
is Floquet isospectral to zero on $G_{d-1}$.
\end{lemma}
\begin{proof}
Set $N=3^{d-1}$, fix all parameters except $t_j$, and write
\[
 H_V(t)=t_jS_j+t_j^{-1}S_j^{-1}+B+\operatorname{diag}(V),
 \qquad
 B=\sum_{\ell\ne j}(t_\ell S_\ell+t_\ell^{-1}S_\ell^{-1}).
\]
Fourier transformation in the $j$th coordinate decomposes the vector space
into three $N$-dimensional eigenspaces of $S_j$, with eigenvalues
$1,\omega,\omega^2$. The compression of $B$ to each of them is the
lower-dimensional free matrix $B$. The compression of
$\operatorname{diag}(V)$ is $\operatorname{diag}(\mathcal A_jV)$, as follows
by taking the normalized inner product in the $j$th coordinate.

Fix $k\in\{0,1,2\}$ and put $\lambda=\omega^kt_j+\mu$. The two other
diagonal blocks of $\lambda I-H_V(t)$ have leading terms
$(\omega^k-\omega^r)t_j I_N$, $r\ne k$, whereas all off-diagonal blocks
are bounded as $t_j\to\infty$. Block elimination, or expansion of the
determinant by its leading power of $t_j$, gives
\begin{align*}
 &\lim_{t_j\to\infty}t_j^{-2N}
     \det\bigl((\omega^kt_j+\mu)I-H_V(t)\bigr)\\
 &\hspace{12mm}=
 \prod_{r\ne k}(\omega^k-\omega^r)^N
 \det\bigl(\mu I-B-\operatorname{diag}(\mathcal A_jV)\bigr).
\end{align*}
The same limit for $V=0$ has the averaged potential absent. The nonzero
constant factors cancel in \eqref{eq:ternary-charidentity}. This proves
the lower-dimensional characteristic-polynomial identity for every $\mu$
and every choice of the other parameters.
\end{proof}

Expand the potential in characters:
\begin{equation}\label{eq:ternary-fourier}
 V(n)=\sum_{a\in G_d}c_a\omega^{a\cdot n},\qquad
 c_a=3^{-d}\sum_{n\in G_d}V(n)\omega^{-a\cdot n}.
\end{equation}
If rigidity has been established in dimension $d-1$,
Lemma~\ref{lem:average-descent} gives $\mathcal A_jV=0$ for every $j$.
In Fourier coordinates this means
\begin{equation}\label{eq:ternary-support}
 c_a=0\quad\hbox{whenever some }a_j=0.
\end{equation}
The only possible remaining frequencies therefore lie in
$\mathcal S_d=\{1,-1\}^d\subset G_d$, and in particular $c_0=0$.
The zero-dimensional initial case is a scalar: equality of its spectrum
with zero forces that scalar to vanish. This supplies the base for the
dimension induction.

\subsection{Exact second- and third-order identities}

We now assume \eqref{eq:ternary-support} and derive the identities that
eliminate the remaining coefficients. Set $t_j=sx_j$ and
$\varepsilon=s^{-1}$, where the $x_j$ are independent nonzero variables.
In the character basis, division of \eqref{eq:ternary-gauge} by $s$ gives
\begin{equation}\label{eq:ternary-perturbation}
 K(\varepsilon)=D+\varepsilon C+\varepsilon^2E,
 \quad D=\operatorname{diag}(L_a),\quad E=\operatorname{diag}(M_a),
 \quad C_{ab}=c_{a-b},
\end{equation}
where
\[
 L_a=\sum_{j=1}^dx_j\omega^{a_j},\qquad
 M_a=\sum_{j=1}^dx_j^{-1}\omega^{-a_j}.
\]
Work over the rational function field $F=\C(x_1,\ldots,x_d)$. The elements
$L_a$ are pairwise distinct: different indices have different coefficient
vectors in their linear forms. Thus the free eigenvalue through $L_0$
is simple, and \eqref{eq:ternary-charidentity} forces the corresponding
formal eigenvalue of $K(\varepsilon)$ to be exactly
\[
 L_0+\varepsilon^2M_0.
\]
There is no ambiguity from a permutation of branches, since their constant
terms are distinct. Define their gaps
\begin{equation}\label{eq:ternary-gaps}
 P_a=L_0-L_a=\sum_{j=1}^d(1-\omega^{a_j})x_j\qquad(a\ne0).
\end{equation}

\begin{lemma}[The two perturbation identities]\label{lem:ternary-identities}
The Fourier coefficients satisfy the following rational identities in $F$:
\begin{align}
 \sum_{a\ne0}\frac{c_ac_{-a}}{P_a}&=0,
       \label{eq:ternary-quadratic}\\
 \sum_{a,b\ne0}
   \frac{c_{-a}c_{a-b}c_b}{P_aP_b}&=0.
       \label{eq:ternary-cubic}
\end{align}
\end{lemma}
\begin{proof}
Choose the formal eigenvector
$u(\varepsilon)=e_0+\varepsilon u^{(1)}+\varepsilon^2u^{(2)}+\cdots$
with zeroth coordinate equal to one. Such a vector exists uniquely by
the simple-eigenvalue equation over $F[[\varepsilon]]$; in particular,
$u_0^{(r)}=0$ for every $r>0$.
The order-one equation from
$K(\varepsilon)u=(L_0+\varepsilon^2M_0)u$ gives
\[
 u_a^{(1)}=\frac{C_{a0}}{P_a}\qquad(a\ne0).
\]
At order two, the nonzero coordinates give
\[
 u_a^{(2)}=
   \sum_{b\ne0}\frac{C_{ab}C_{b0}}{P_aP_b}.
\]
Here $Ee_0$ has no nonzero coordinate. The zeroth coordinate at the same
order is
$\sum_{a\ne0}C_{0a}u_a^{(1)}+M_0=M_0$, which is
\eqref{eq:ternary-quadratic}. At order three, the zeroth coordinate gives
$\sum_{a\ne0}C_{0a}u_a^{(2)}=0$, yielding
\eqref{eq:ternary-cubic}. Both $Eu^{(1)}$ and $M_0u^{(1)}$ have zero
zeroth coordinate. This accounts for all possible contributions from the
$\varepsilon^2E$ term; there is no omitted mixed term.
\end{proof}

By \eqref{eq:ternary-support}, the nonzero terms in the quadratic identity
have $a\in\mathcal S_d$. In a nonzero term of the cubic identity the three
steps $-a$, $a-b$, and $b$ all lie in $\mathcal S_d$. In each coordinate
they sum to zero in $\Z/3\Z$. Three signs in $\{1,-1\}$ sum to zero
modulo three only when all three are equal. It follows that $b=-a$, and
the cubic identity becomes
\begin{equation}\label{eq:ternary-cubic-pairs}
 \sum_{\{a,-a\}\subset\mathcal S_d}
     \frac{c_a^3+c_{-a}^3}{P_aP_{-a}}=0,
\end{equation}
where the sum takes one representative of each antipodal pair.

\subsection{Separation by the gap hyperplanes}

Put $A=1-\omega$, $B=1-\omega^2$, and $S=x_1+\cdots+x_d$. Then
\[
 A+B=AB=3,\qquad A\ne B,\qquad A^2\ne B^2.
\]
For a sign vector $a$, the coefficients of $P_a$ are $A$ and $B$.
Call $a$ mixed if both signs occur.

\begin{lemma}[Distinct mixed gap hyperplanes]\label{lem:mixed-gaps}
If $a\in\mathcal S_d$ is mixed, $P_a$ is not proportional to $P_b$ for
any distinct $b\in\mathcal S_d$. The only proportional pair of gap forms
is that for the two constant sign vectors, with forms $AS$ and $BS$.
\end{lemma}
\begin{proof}
Suppose $P_b=rP_a$. Since both coefficient values occur in $P_a$, the
nonzero scalar $r$ must map the two-element set $\{A,B\}$ to itself.
If it fixes both elements then $r=1$ and $b=a$. If it interchanges them,
then $rA=B$ and $rB=A$, contradicting $A^2\ne B^2$. A form with two
different coefficient values cannot be proportional to a constant-coefficient
form. The assertion for the two constant vectors is immediate.
\end{proof}

For mixed $a$, multiply \eqref{eq:ternary-quadratic} by $P_a$ and reduce
modulo the prime ideal $(P_a)$ in $\C[x_1,\ldots,x_d]$, after localizing
the other denominator factors. By Lemma~\ref{lem:mixed-gaps}, only its
single term survives. Hence $c_ac_{-a}=0$ for every mixed $a$.
Let $c_+$ and $c_-$ denote the coefficients of the two constant sign
vectors. The remaining quadratic identity is
\[
 \frac{c_+c_-}{S}\left(\frac1A+\frac1B\right)
      =\frac{c_+c_-}{S}=0.
\]
We have therefore established
\begin{equation}\label{eq:ternary-products}
 c_ac_{-a}=0\qquad(a\in\mathcal S_d).
\end{equation}

Apply the same operation to \eqref{eq:ternary-cubic-pairs}. For mixed
$a$, its denominator factor $P_a$ occurs in only that pair, and $P_{-a}$
is nonzero modulo $(P_a)$. This gives $c_a^3+c_{-a}^3=0$.
The remaining constant pair has denominator $AB S^2$, so its numerator
also vanishes. Thus
\begin{equation}\label{eq:ternary-cubes}
 c_a^3+c_{-a}^3=0\qquad(a\in\mathcal S_d).
\end{equation}
Equations \eqref{eq:ternary-products} and \eqref{eq:ternary-cubes} imply
$c_a=c_{-a}=0$: if one were nonzero, the product equation would force
the other to be zero, contradicting the cube equation. In dimension one
there are no mixed vectors, and the constant-pair calculations above give
the entire argument without any hyperplane restriction.

\begin{proof}[Proof of Theorem~\ref{thm:ternary-rigidity}]
We induct on the dimension, beginning with the scalar case $d=0$.
For a zero-isospectral potential in dimension $d$, the induction hypothesis
and Lemma~\ref{lem:average-descent} give
\eqref{eq:ternary-support}. Lemma~\ref{lem:ternary-identities} and the
gap-hyperplane argument then show that every remaining Fourier coefficient
is zero. Fourier inversion proves $V\equiv0$. This establishes rigidity
for every $(3\Z)^d$.

For the rectangular classification, \cite[Proposition~4.2]{CFK} states
the following exact equivalence: when the last period is one or two, a
nonzero zero-isospectral potential exists in that dimension if and only
if one exists after deleting that coordinate. Repeatedly apply this
equivalence to a period vector with all $q_j\leq3$. The remaining vector
has all entries equal to three, where rigidity has just been proved, or
has dimension zero. Thus no nonzero potential exists.
Conversely, if some $q_j>3$, existence is
\cite[Theorem~C]{CFK}. These two directions prove
\eqref{eq:rectangular-threshold}.
\end{proof}

The proof uses identities in independent Floquet variables. It does not
infer rigidity from a single fixed Floquet matrix or impose reality,
separability, or smaller least period on the potential. The explicitly
open family in \cite[Section~4.4]{CFK} is settled in full by the
period-three assertion. The final if-and-only-if statement completes the
rectangular classification by adjoining the cited prior results. General
oblique lattices remain outside this conclusion.

\section{T80: optimal quadratic stabilization in equivariant Ehrhart theory}
\label{sec:ehrhart-stabilization}

The target in this section is Conjecture~8.5 of Clarke
\cite{Clarke26}. It is a coefficientwise inequality for the equivariant
Ehrhart series of an arbitrary lattice polytope with an involution. We first
formulate and solve a parameter-construction problem, then apply its answer
to that conjecture without a polynomiality or integer-decomposition
assumption on the polytope.

\subsection{The original inequality and the auxiliary problem}

Let $L$ be a lattice of rank $n+1$, equipped with a primitive height
functional. Let $P$ be a full-dimensional lattice polytope in the affine
hyperplane of height one, preserved by a nontrivial involution $\sigma$ of
$L$. Write $G=\{1,\sigma\}$ and let $\varepsilon$ denote its sign character.
Let $p$ and $q$ be the dimensions of the $+1$ and $-1$ eigenspaces of the
ambient representation. Thus
\[
 p+q=n+1,\qquad p\geq1,\qquad q\geq1.
\]
The restriction $q\geq1$ excludes the trivial action; it is necessary for
the universal optimality statement below.

Let $R=\mathbb C[\operatorname{cone}(P)\cap L]$ be the Ehrhart ring, graded
by height. Its equivariant Hilbert series is
\[
 E_P(t)=\sum_{j\geq0}[R_j]t^j
       =\frac{\varphi_P(t)}{D(t)},\qquad
 D(t)=(1-t)^p(1-\varepsilon t)^q,
\]
in the representation ring $R(G)=\mathbb Z[\varepsilon]/(\varepsilon^2-1)$.
An element of this ring is \emph{effective} if it is the character of an
actual representation, equivalently if both its trivial and sign
multiplicities are nonnegative. The same term is used coefficientwise for
a power series.

The evaluation $\varphi_P(t)(1)$ is the ordinary polynomial $h^*_P(t)$.
The fixed polytope has denominator dividing two. Write its Ehrhart series
as $H(t)/(1-t^2)^p$, where $H$ has nonnegative integer coefficients. Then
\begin{equation}\label{eq:ehrhart-fixed-numerator}
 \varphi_P(t)(\sigma)=(1+t)^{q-p}H(t).
\end{equation}
These are the conventions of \cite[Section~2 and Lemma~8.1]{Clarke26}.
Clarke's Conjecture~8.5 asserts
\begin{equation}\label{eq:clarke-target}
 \begin{cases}
 (1+t)^{p-q}h^*_P(t)\ \geq\ H(t),&p\geq q,\\
 h^*_P(t)\ \geq\ (1+t)^{q-p}H(t),&p<q,
 \end{cases}
\end{equation}
where each inequality compares coefficients.

Our auxiliary question is the following: how many invariant quadratic
parameters must replace invariant linear parameters to obtain a
homogeneous system of parameters of $R$, while keeping exactly $q$
anti-invariant linear parameters? A second question asks for the least
exponent $s=s(p,q)$ such that $(1+t)^s\varphi_P(t)$ is effective for
\emph{every} polytope of ambient type $(p,q)$. The first answer depends on
a base locus; the universal answer is $\max(p-q,0)$.

Put $X=\operatorname{Proj}R$, so $\dim X=p+q-1$. Write $W_+$ and $W_-$
for the invariant and anti-invariant subspaces of $R_1$, and use the
base-locus convention
\[
 Z_+=V(W_-),\qquad Z_-=V(W_+),\qquad b=\dim Z_-.
\]
If $Z_-$ is empty, set $b=-\infty$ and interpret
$\max(0,b-q+1)$ as zero.

\subsection{Base loci and generic sections}

\begin{lemma}\label{lem:ehrhart-base-loci}
The space $R_1$ is basepoint-free on $X$, and
\[
 Z_+\cap Z_-=\varnothing,\qquad
 \dim Z_+=p-1,\qquad \dim Z_-\leq p-1.
\]
The invariant subspace $(R_2)^G$ is also basepoint-free.
\end{lemma}

\begin{proof}
The linear base-locus assertions are
\cite[Lemmas~7.1(a) and~7.3]{Clarke26}. We recall their geometric content
to make clear that no assumption that $R$ is generated in degree one is
being made. Every point of the projective toric variety $X$ lies on the
orbit associated to a nonempty face $F$ of $P$. A monomial corresponding
to any vertex of $F$ is nonzero at that point. Hence $R_1$ has no base
point, and the two eigenspace base loci are disjoint.

On a face orbit, the equations defining either $Z_+$ or $Z_-$ force the
face to be $\sigma$-invariant: otherwise a vertex monomial supported on
$F$ is paired with a monomial vanishing on that orbit. For an invariant
face, the equations on its torus characters have the form
$\chi(\sigma u)=\chi(u)$ or $\chi(\sigma u)=-\chi(u)$ for its lattice
points $u$. Their solution set is either empty or a coset of a torus
whose dimension, after quotienting by projective scaling, is
$\dim F^\sigma\leq p-1$. For the invariant equations on the full face
$P$, the trivial character supplies a solution and the dimension is
$p-1$. Taking the finite union over face orbits gives the asserted bounds.

Finally, at any point of $X$, at least one eigenlinear form
$f\in W_+\cup W_-$ is nonzero. Its square $f^2$ is an invariant element
of degree two, still nonzero there. This proves the last assertion.
\end{proof}

\begin{lemma}\label{lem:ehrhart-mixed-incidence}
Let $0\leq k\leq p$. A homogeneous system of parameters consisting of
$p-k$ invariant linear forms, $q$ anti-invariant linear forms, and $k$
invariant quadratic forms exists if and only if
\begin{equation}\label{eq:ehrhart-budget}
 k\geq\delta(P):=\max(0,b-q+1).
\end{equation}
\end{lemma}

\begin{proof}
First suppose \eqref{eq:ehrhart-budget} holds. Choose the linear forms
from the affine parameter space
\[
 \mathcal U=W_+^{\,p-k}\times W_-^{\,q}.
\]
Over $X\setminus(Z_+\cup Z_-)$, the condition that all chosen forms vanish
at a specified point imposes $p-k+q$ independent linear conditions on
$\mathcal U$: each evaluation acts nontrivially in its own tuple
coordinate. Over $Z_+$, it imposes $p-k$ conditions; over $Z_-$, it
imposes $q$. Thus the corresponding incidence strata have dimensions at
most, respectively,
\begin{align*}
 \dim\mathcal U+(p+q-1)-(p-k+q)&=\dim\mathcal U+k-1,\\
 \dim\mathcal U+(p-1)-(p-k)&=\dim\mathcal U+k-1,\\
 \dim\mathcal U+b-q&\leq\dim\mathcal U+k-1.
\end{align*}
The generic-fiber dimension theorem now implies that a general tuple of
linear forms has common zero locus $Y\subset X$ of dimension at most
$k-1$. A stratum not dominating $\mathcal U$ has empty generic fiber and
does not affect this conclusion. In particular, if $k=0$, then $Y$ is
empty.

By Lemma~\ref{lem:ehrhart-base-loci}, invariant quadratics are
basepoint-free. Therefore $k$ general invariant quadratic sections have
no common zero on $Y$: successively they cut positive-dimensional
components by one dimension and avoid any remaining finite set of
points. The full collection has no common projective zero. By the graded
Nullstellensatz, its homogeneous ideal has radical equal to the
irrelevant ideal, and the quotient ring is finite-dimensional. Since
there are $p+q=\dim R$ forms, they are a homogeneous system of parameters.

Conversely, consider any system with the stated degree and character
pattern, and suppose $Z_-$ is nonempty. On a component of $Z_-$ of
dimension $b$, all invariant linear forms vanish identically. The
remaining $q+k$ positive-degree sections must therefore have no common
zero on that projective component. This requires $q+k\geq b+1$: at most
$b$ positive-degree hypersurfaces cannot have empty intersection on a
projective variety of dimension $b$. Equivalently, its homogeneous
coordinate ring has dimension $b+1$, and Krull's height theorem prevents
an ideal generated by at most $b$ positive-degree elements from having a
finite quotient. Hence $k\geq b-q+1$. Together with $k\geq0$, this proves
necessity, including the empty-base-locus case.
\end{proof}

\subsection{The stabilization theorem and its application}

\begin{theorem}[Optimal quadratic stabilization and Clarke's inequality]
\label{thm:optimal-ehrhart-stabilization}
For the order-two lattice action above, set
\[
 \delta(P)=\max(0,\dim V(W_+)-q+1),\qquad r=\max(p-q,0).
\]
Then $0\leq\delta(P)\leq r$, and there is a finite-dimensional graded
$G$-algebra with character
\begin{equation}\label{eq:ehrhart-artinian-character}
 (1+t)^{\delta(P)}\varphi_P(t).
\end{equation}
The number $\delta(P)$ is the minimum number of invariant quadratics in
the parameter pattern of Lemma~\ref{lem:ehrhart-mixed-incidence}.

For each fixed pair $p,q\geq1$, the least integer $s\geq0$ for which
$(1+t)^s\varphi_P(t)$ is effective for every polytope of ambient type
$(p,q)$ is exactly $r$. In particular, the original coefficient
inequalities \eqref{eq:clarke-target}, namely Conjecture~8.5 of
\cite{Clarke26}, hold for every such polytope.
\end{theorem}

\begin{proof}
The bound on $\delta(P)$ follows from
Lemma~\ref{lem:ehrhart-base-loci}. Apply
Lemma~\ref{lem:ehrhart-mixed-incidence} with $k=\delta(P)$. The Ehrhart
ring is a normal affine semigroup ring and is Cohen--Macaulay by
Hochster's theorem \cite{Hochster72}. Consequently its homogeneous
system of parameters is a regular sequence. Taking its equivariant
Koszul complex gives the character of the Artinian quotient as
\begin{align*}
 E_P(t)(1-t)^{p-\delta(P)}(1-\varepsilon t)^q
                 (1-t^2)^{\delta(P)}
 &=\varphi_P(t)(1+t)^{\delta(P)}.
\end{align*}
This is an actual graded representation, proving
\eqref{eq:ehrhart-artinian-character}. The parameter minimality is
precisely the necessity assertion of
Lemma~\ref{lem:ehrhart-mixed-incidence}.

Since $r\geq\delta(P)$, multiplication by the trivially acted-on
polynomial $(1+t)^{r-\delta(P)}$ preserves effectiveness. Thus
$(1+t)^r\varphi_P(t)$ is always effective. Its evaluations at the
identity and at $\sigma$ are
\[
 (1+t)^r h^*_P(t),\qquad (1+t)^{r+q-p}H(t),
\]
respectively. For an effective character of a group of order two, its
identity value is at least the absolute value of its value at $\sigma$.
The second displayed polynomial has nonnegative coefficients. Comparing
each degree therefore proves exactly \eqref{eq:clarke-target} in both
cases $p\geq q$ and $p<q$.

It remains to prove the sharpness of the universal exponent. There is
nothing to prove when $p\leq q$, since then $r=0$. Suppose $r=p-q>0$.
Start with the cube $C=[0,1]^{r+1}$ and reflect its first coordinate:
\[
 (x_1,x_2,\ldots,x_{r+1})\longmapsto
       (1-x_1,x_2,\ldots,x_{r+1}).
\]
On the homogeneous lattice $L_0=\mathbb Z^{r+2}$ with coordinates
$(x_1,\ldots,x_{r+1},h)$, this is the integral involution
\[
 (x_1,x_2,\ldots,x_{r+1},h)\longmapsto
       (h-x_1,x_2,\ldots,x_{r+1},h).
\]
Its ambient type is $(r+1,1)$. The fixed polytope has $x_1=1/2$, so its
lattice-point count in dilation $j$ is $(j+1)^r$ for even $j$, and zero
for odd $j$. Its Ehrhart series is therefore
\[
 E_{C^\sigma}(t)=\sum_{\ell\geq0}(2\ell+1)^r t^{2\ell}.
\]
The generating series of the degree-$r$ polynomial $(2\ell+1)^r$ has a
genuine pole of order $r+1$ at $t^2=1$: its leading contribution is
$2^r r!/(1-t^2)^{r+1}$. In particular, the pole at $t=-1$ has order
$r+1$. Multiplication by
$D_C(t)(\sigma)=(1-t)^{r+1}(1+t)$ leaves a pole of order exactly $r$ in
$\varphi_C(t)(\sigma)$.

To obtain arbitrary $q\geq1$ while retaining this pole, take the standard
lattice join with $q-1$ unimodular segments, each with the involution
swapping its endpoints. More explicitly, use
\[
 L=L_0\oplus(\mathbb Z^2)^{q-1},\qquad
 \operatorname{ht}(x,h;y)=h+\sum_{a=1}^{q-1}(y_{a,0}+y_{a,1}),
\]
and take the convex hull of the embedded $(C,1)$ and the embedded
segments $\operatorname{conv}(e_{a,0},e_{a,1})$ in height one. Apply the
cube involution on $L_0$ and swap each added coordinate pair. The cone
over this join is the product of the cube cone with $q-1$ nonnegative
quadrants. Its lattice semigroup is consequently a direct product, and
its Ehrhart ring is the graded tensor product of $R_C$ with $q-1$
two-variable polynomial rings having swapped variables.

Each variable pair contributes ambient representation $1+\varepsilon$
and Hilbert series $1/((1-t)(1-\varepsilon t))$. Its determinant factor
cancels, so the join has equivariant numerator exactly $\varphi_C(t)$.
Its ambient type is $(r+q,q)=(p,q)$. Therefore a multiplier $(1+t)^s$
with $s<r$ leaves a pole in its evaluation at $\sigma$.

An effective series whose identity evaluation is a polynomial must itself
be a polynomial: in every degree beyond the degree of that evaluation,
an effective representation of dimension zero is zero. The identity
evaluation of $(1+t)^s\varphi_P(t)$ is the polynomial
$(1+t)^s h^*_P(t)$. Hence the remaining pole precludes effectiveness for
$s<r$. This proves the universal lower bound and completes the theorem.
\end{proof}

The two optimality assertions must be distinguished. For a particular
polytope, $\delta(P)$ is optimal for the specified degree and character
pattern of parameters. It need not be the least exponent making its
series effective: \cite[Example~5.3]{Clarke26} already exhibits an
effective series with an obstruction to linear parameters. The exponent
$r$ is instead optimal after quantifying over all polytopes of fixed
ambient type. The construction above proves that universal statement.

The base-locus description, generic incidence argument, and Koszul
character formula are prior tools. The contribution proposed here is
the exact mixed-degree parameter budget, its sharp universal
stabilization consequence, and the resulting proof of the original
coefficient inequality. This distinction concerns historical attribution;
it does not add an unproved hypothesis to the mathematical argument.

\section{Attribution, novelty, and the requested success conditions}
\label{sec:historical-audit}

The prior problem statements are pinned to specific primary sources:
\cite[Conjecture~4]{KMp}, \cite[Conjecture~12]{KM3},
\cite[Section~4.4]{CFK}, and \cite[Conjecture~8.5]{Clarke26}.
The first two sources originate in 2015 and 2013; the latter two in 2026.
The spectral target is explicitly open in the cited source, but is not
numbered or given a standalone conjecture name there. It must not be
relabelled as a previously named conjecture. The nested binary special case
\cite[Conjecture~4]{KKM2} is not counted as another independent target.

The arithmetic construction uses the classical congruence
$B(z)^p\equiv B(z^p)\pmod p$. Its specific contribution here is the
substitution-compatible correction term, which produces the required
valuation at every level and a full kernel description. The Cantor proof
returns every coefficient to the prescribed ring before claiming the
original conjecture. No coefficient-ring enlargement is hidden in the
final statement.

The spectral argument uses earlier Fourier and rational-pole techniques
\cite{Liu23}. Its additional period-three step is that all surviving
three-step Fourier cycles follow a single direction. The resulting cubic
relations remove the remaining complex coefficients without a positivity
assumption. The contribution is this specific closure argument, not a claim
to have invented perturbation theory or pole extraction.

The Ehrhart construction extends Clarke's base-locus and generic-section
framework \cite{Clarke26}. The exact quadratic budget and its character
factor are the additional ingredients. The Cohen--Macaulay and equivariant
Koszul facts are prior \cite{Hochster72,Stapledon24}. A different proof of
Clarke's dilation Conjecture~8.7 was also examined, but its central
determinant-ratio and coefficient-extraction mechanism is already explicit
in \cite[Section~2.5, equation~(7), and Lemma~2.33]{Stapledon24}. It is not
included as an additional new-method success in this manuscript.

\paragraph{Correction to the preceding T71--T75 attribution.}
The even-period Fourier-square quadratic structure used in the preceding
T71 draft is already explicit in Korotyaev and Kutsenko
\cite[Theorem~2.4, Lemma~2.6, and equation~(2.31)]{KK05}.
That structure must be credited to their 2005 preprint and subsequent
publication. The correction narrows the earlier novelty description; it
does not itself contradict the algebraic consequences proved in that draft.
The present T79 argument is a separate multidimensional rigidity argument,
and no part of T76--T80 counts that older even-period structure as a new
principle.

\paragraph{What has and has not been established.}
The displayed arguments address the stated source problems unchanged and
without conjectural assumptions. Searches through 24 October 2026 examined
the current source versions, relevant author publication lists, and topic
and conjecture-number searches using two search services. No identical
earlier resolution was located. Such a search cannot prove universal
absence from the literature. Nor does a new construction using classical
tools establish that its general principle is historically unprecedented.
Consequently this manuscript does \emph{not} certify all six requested
historical success conditions. It supplies the complete mathematical
arguments and an auditable, bounded priority assessment. The numbering
T76--T80 organizes five theorems, with three numbered prior conjectures and
one prior unnumbered problem as applications.

\appendix
\section{Exact checks and reproducibility}

The accompanying verification bundle contains executable Python scripts and
their output. All checks use exact integer or rational arithmetic. They
support the proofs but establish only the finite cases actually executed.

\begin{center}\small
\begin{tabular}{@{}p{.31\textwidth}p{.60\textwidth}@{}}
\toprule
Check & Executed range and result \\
\midrule
Additive Frobenius tower & First 256 coefficients; $p=2$ through level 4,
$p=3$ through level 3, and $p=5$ through level 2. Every predicted
coefficientwise divisibility holds.\\
Cantor tower & First 256 coefficients at levels 0, 1, 2, with degrees
$2,6,18$ and valuations at least $1,4,13$. The displayed original-ring
degree-six relation modulo $81$ also passes.\\
Spectral pole separation & Dimensions $1$ through $5$: the evaluated rational
bases have exact ranks $1,2,4,8,16$, and every admissible three-step cycle
uses a single Fourier direction.\\
Reflection cubes & Dimensions $2$ through $7$: exact Ehrhart and fixed-point
numerators give nonnegative integral character multiplicities after the
predicted stabilization; the pole order at $-1$ is the claimed one.\\
\bottomrule
\end{tabular}
\end{center}

The full proofs were also checked by independent internal derivations:
coefficient isolation and valuation descent for T76--T78; formal eigenvector
recursion and hyperplane residues for T79; and incidence dimensions,
non-degree-one generation, and equivariant Koszul factors for T80.
These checks should not be described as external peer review or formal
machine verification.

The source is standalone: compiling it twice with \texttt{pdflatex} requires
no separate figures or bibliography files. The verification bundle includes
its own run instructions and records the exact finite scope of each script.


% Primary-source verified bibliography, audited 24 October 2026.
% The key Stapledon24 is retained for compatibility; the cited current version is v3 (2025).
% KM3 remains a preprint in the authors' current publication list; no journal is asserted.
\begin{thebibliography}{99}

\bibitem{KMp}
C.~Krattenthaler and T.~W.~M\"uller,
\emph{A method for determining the mod-$p^k$ behaviour of recursive sequences},
Aequationes Mathematicae \textbf{99} (2025), 2505--2547.
\href{https://doi.org/10.1007/s00010-025-01195-0}{doi:10.1007/s00010-025-01195-0}.
First preprint: 2015; \href{https://arxiv.org/abs/1508.02580}{arXiv:1508.02580}, v2, 28 July 2025.

\bibitem{KM3}
C.~Krattenthaler and T.~W.~M\"uller,
\emph{A method for determining the mod-$3^k$ behaviour of recursive sequences},
preprint (2013), \href{https://arxiv.org/abs/1308.2856}{arXiv:1308.2856v1}, 13 August 2013.

\bibitem{KKM2}
M.~Kauers, C.~Krattenthaler and T.~W.~M\"uller,
\emph{A method for determining the mod-$2^k$ behaviour of recursive sequences, with applications to subgroup counting},
Electronic Journal of Combinatorics \textbf{18}(2) (2012), Paper P37.
\href{https://doi.org/10.37236/2433}{doi:10.37236/2433};
\href{https://arxiv.org/abs/1107.2015}{arXiv:1107.2015} (first submitted 2011).

\bibitem{CFK}
J.~Cobb, M.~Faust and A.~Kretschmer,
\emph{Inverse eigenvalue problems, Floquet isospectrality and the Hilbert--Chow morphism},
preprint (2026), \href{https://arxiv.org/abs/2601.09604}{arXiv:2601.09604v1}, 14 January 2026.

\bibitem{Liu23}
W.~Liu,
\emph{Floquet isospectrality for periodic graph operators},
Journal of Differential Equations \textbf{374} (2023), 642--653.
\href{https://doi.org/10.1016/j.jde.2023.08.009}{doi:10.1016/j.jde.2023.08.009};
\href{https://arxiv.org/abs/2302.13103}{arXiv:2302.13103}.

\bibitem{Clarke26}
O.~Clarke,
\emph{The equivariant Ehrhart theory of polytopes with order-two symmetries: Part II},
preprint (2026), \href{https://arxiv.org/abs/2610.17988}{arXiv:2610.17988v1}, 13 October 2026.

\bibitem{Stapledon24}
A.~Stapledon,
\emph{Equivariant Ehrhart theory, commutative algebra and invariant triangulations of polytopes},
preprint (2023; revised 2025),
\href{https://arxiv.org/abs/2311.17273}{arXiv:2311.17273v3}, 24 September 2025.

\bibitem{Hochster72}
M.~Hochster,
\emph{Rings of invariants of tori, Cohen--Macaulay rings generated by monomials, and polytopes},
Annals of Mathematics (2) \textbf{96}(2) (1972), 318--337.
\href{https://doi.org/10.2307/1970791}{doi:10.2307/1970791}.

\bibitem{KK05}
E.~Korotyaev and A.~Kutsenko,
\emph{Inverse problem for the discrete 1D Schr\"odinger operator with small periodic potentials},
Communications in Mathematical Physics \textbf{261} (2006), 673--692
(published online 20 September 2005).
\href{https://doi.org/10.1007/s00220-005-1429-z}{doi:10.1007/s00220-005-1429-z};
\href{https://arxiv.org/abs/math/0507297}{arXiv:math/0507297}.

\end{thebibliography}

\end{document}
